% part: intuitionistic-logic
% chapter: semantics
% section: introduction

\documentclass[../../../include/open-logic-section]{subfiles}

\begin{document}

\olfileid{int}{sem}{int}

\olsection{Pambuka}

Ora ana logika kang bisa diandharake kanthi marem tanpa semantik,
lan logika intuisionistik uga mangkono. Ing logika klasik, semantik
adhedhasar !!{valuation}s iku kanonik, dene ana sawetara semantik
kang saingan kanggo logika intuisionistik. Ora ana siji-sijia kang
marem kanthi sampurna ing pangertene menehi andharan teges konektif
kang bisa ditampa sacara intuisionistik.

Semantik adhedhasar !!{relational model}s, kang padha karo semantik
kanggo logika modal, mbokmenawa paling kerep digunakake. Ing semantik
iki, !!{propositional variable}s ditetepake ing donya-donya, lan donya
iku gegandhengan liwat relasi aksesibilitas. Relasi iku tansah urutan
parsial, yaiku refleksif, antisimetris, lan transitif.

Sacara intuitif, donya iki bisa dianggep kahanan kawruh utawa
``kahanan bukti.'' Kahanan~$w'$ bisa diakses saka~$w$ yen lan mung
yen, miturut kabeh kang kita weruhi, $w'$ iku kahanan kawruh
(mbesuk) kang bisa kelakon, yaiku kang cocog karo apa kang diweruhi
ing~$w$. Yen sawijining proposisi wis diweruhi, ora bisa dadi ora
diweruhi maneh. Tegese, saben $!A$ diweruhi ing~$w$ lan~$Rww'$,
$!A$ uga diweruhi ing~$w'$. Dadi ``kawruh'' iku monoton tumrap
relasi aksesibilitas.

Yen ``$!A$ diweruhi'' ditetepake kaya ing logika epistemik minangka
``bener ing kabeh alternatif epistemik,'' mula $!A \land !B$
diweruhi ing~$w$ yen ing kabeh alternatif epistemik, $!A$ lan~$!B$
loro-lorone diweruhi. Nanging amarga kawruh iku monoton lan $R$
refleksif, tegese $!A \land !B$ diweruhi ing $w$ yen lan mung yen
$!A$ lan $!B$ diweruhi ing~$w$. Kanthi alasan kang padha,
$!A \lor !B$ diweruhi ing $w$ yen lan mung yen paling ora salah
sijine diweruhi. Dadi kanggo $\land$ lan $\lor$, syarat kabeneran
konektif padha karo ing logika klasik.

Nanging, syarat kabeneran kondisional beda karo ing logika klasik.
$!A \lif !B$ diweruhi ing~$w$ yen lan mung yen ora ana $w'$ kanthi
$Rww'$ kang $!A$ diweruhi nanging $!B$ ora diweruhi. Iki ora padha
karo syarat yen $!A$ ora diweruhi utawa $!B$~diweruhi ing~$w$.
Amarga yen $!A$ lan $!B$ loro-lorone ora kita weruhi ing~$w$, bisa wae ana
kahanan epistemik mbesuk~$w'$ kanthi $Rww'$ kang ing $w'$, $!A$
diweruhi nanging~$!B$ tetep ora diweruhi.

Kita weruh $\lnot !A$ mung yen ora ana kahanan epistemik mbesuk kang
bisa kelakon lan ing kono kita weruh~$!A$. Gagasane yaiku yen $!A$
bisa diweruhi, mula ing sawenehing kahanan epistemik mbesuk kang bisa
kelakon,~$!A$ dadi diweruhi. Amarga kita ora bisa weruh $\lfalse$,
ing kahanan epistemik mbesuk iku kita bakal weruh $!A$ nanging ora
weruh~$\lfalse$.

Ing interpretasi iki, prinsip Hukum Tengah kang Disingkirake ora
lumaku. Amarga ana sawenehing $!A$ kang durung kita weruhi nanging
bisa dadi diweruhi mengko. Kanggo !!a{formula}~$!A$ kaya mangkono,
$!A$ lan~$\lnot !A$ loro-lorone ora diweruhi, mula
$!A \lor \lnot !A$ ora diweruhi. Nanging kita weruh, upamane,
yen $\lnot(!A \land \lnot !A)$. Amarga ora bisa ana kahanan mbesuk
kang kita weruh $!A$ lan $\lnot !A$ bebarengan, lan kita weruh iki
tanpa gumantung marang apa kita weruh~$!A$ utawa $\lnot !A$ utawa ora.

!!^{relational model}s dudu siji-sijine semantik kang kasedhiya
kanggo logika intuisionistik. Semantik topologis iku pilihan liyane:
ing kene proposisi diinterpretasikake minangka himpunan kabuka ing
ruang topologis, lan konektif diinterpretasikake minangka operasi
ing himpunan iku (upamane, $\land$ cocog karo irisan).

\end{document}
